% ─────────────────────────────────────────────────────────────────────────────
% report.tex — Báo cáo tóm tắt kiến trúc hệ thống gợi ý H&M
% Compile: pdflatex report.tex (hoặc xelatex report.tex)
% ─────────────────────────────────────────────────────────────────────────────
\documentclass[12pt,a4paper]{article}

% ── Packages ────────────────────────────────────────────────────────────────
\usepackage[utf8]{inputenc}
\usepackage[vietnamese]{babel}
\usepackage[T5]{fontenc}
\usepackage{geometry}
\geometry{margin=2.5cm}

\usepackage{amsmath, amssymb}
\usepackage{graphicx}
\usepackage{booktabs}
\usepackage{enumitem}
\usepackage{hyperref}
\usepackage{xcolor}
\usepackage{listings}
\usepackage{caption}
\usepackage{float}

\hypersetup{
    colorlinks=true,
    linkcolor=blue!70!black,
    citecolor=green!60!black,
    urlcolor=blue!60!black,
}

\lstset{
    language=Python,
    basicstyle=\ttfamily\small,
    keywordstyle=\color{blue!80!black}\bfseries,
    commentstyle=\color{green!50!black}\itshape,
    stringstyle=\color{red!60!black},
    breaklines=true,
    frame=single,
    rulecolor=\color{gray!30},
    backgroundcolor=\color{gray!5},
    numbers=left,
    numberstyle=\tiny\color{gray},
    tabsize=4,
}

% ── Title ───────────────────────────────────────────────────────────────────
\title{%
    \textbf{Báo cáo kỹ thuật}\\[6pt]
    \Large Hệ thống Gợi ý Thời trang H\&M\\
    \large Hybrid Neural Collaborative Filtering kết hợp Visual Features
}
\author{Nhóm phát triển — Intro AI}
\date{\today}

% ─────────────────────────────────────────────────────────────────────────────
\begin{document}
\maketitle
\tableofcontents
\newpage

% ═════════════════════════════════════════════════════════════════════════════
\section{Tổng quan dự án}
% ═════════════════════════════════════════════════════════════════════════════

Dự án xây dựng một hệ thống gợi ý sản phẩm thời trang (Fashion Recommendation System)
dựa trên bộ dữ liệu \textbf{H\&M Personalized Fashion Recommendations} từ Kaggle.
Hệ thống kết hợp kỹ thuật \textit{Neural Collaborative Filtering (NCF)} với
\textit{Visual Feature Extraction} để tạo ra một mô hình Hybrid có khả năng
cá nhân hoá gợi ý cho từng khách hàng.

\subsection{Kiến trúc tổng thể}

Hệ thống gồm 5 module chính:

\begin{table}[H]
\centering
\caption{Danh sách các module và chức năng}
\begin{tabular}{lll}
\toprule
\textbf{Module} & \textbf{Tệp nguồn} & \textbf{Vai trò} \\
\midrule
Data Loader      & \texttt{data\_loader.py}   & Đọc dữ liệu, tiền xử lý, tạo Dataset \\
Models           & \texttt{models.py}         & NCF, FeatureExtractor, HybridModel \\
Metrics          & \texttt{metrics.py}        & MAP@12, NDCG@K, Hit Rate \\
Training         & \texttt{train\_hybrid.py}  & Vòng lặp huấn luyện, Negative Sampling \\
Demo UI          & \texttt{app.py}            & Giao diện Streamlit \\
\bottomrule
\end{tabular}
\end{table}

% ═════════════════════════════════════════════════════════════════════════════
\section{Module \texttt{data\_loader.py} — Quản lý dữ liệu}
% ═════════════════════════════════════════════════════════════════════════════

Module này đảm nhiệm toàn bộ việc đọc và tiền xử lý dữ liệu, sử dụng thư viện
\textbf{Polars} để tối ưu hiệu năng so với Pandas truyền thống.

\subsection{Class \texttt{DataLoaderPolars}}

\begin{itemize}[leftmargin=2em]
    \item Đọc 4 tệp Parquet: \texttt{articles\_cleaned}, \texttt{customers\_fixed},
          \texttt{hm\_train}, \texttt{hm\_test}.
    \item Phương thức \texttt{join\_features(df)} thực hiện \texttt{LEFT JOIN}
          giữa bảng giao dịch với bảng khách hàng (qua \texttt{user\_id}) và
          bảng sản phẩm (qua \texttt{item\_id} hoặc \texttt{article\_id}).
\end{itemize}

\subsection{Class \texttt{PopularityBaseline}}

Baseline đơn giản nhất: đếm tần suất mua của từng \texttt{item\_id} trong tập
huấn luyện, lấy 12 sản phẩm bán chạy nhất làm gợi ý mặc định cho mọi người dùng.

\subsection{Class \texttt{MatrixFactorization}}

Mô hình phân rã ma trận (Matrix Factorization) bằng PyTorch:
\begin{equation}
    \hat{r}_{u,i} = \mathbf{p}_u^{\top} \mathbf{q}_i + b_u + b_i
\end{equation}
trong đó $\mathbf{p}_u$, $\mathbf{q}_i$ là embedding vectors,
$b_u$, $b_i$ là bias tương ứng.

\subsection{Class \texttt{HMDataset} và hàm \texttt{calculate\_map12}}

\begin{itemize}[leftmargin=2em]
    \item \texttt{HMDataset}: Wrapper cho PyTorch \texttt{Dataset},
          chuyển đổi các mảng \texttt{user\_id}, \texttt{item\_id}, \texttt{label}
          thành tensors.
    \item \texttt{calculate\_map12}: Tính Mean Average Precision @ 12
          trên tập test.
\end{itemize}

% ═════════════════════════════════════════════════════════════════════════════
\section{Module \texttt{models.py} — Kiến trúc mô hình}
% ═════════════════════════════════════════════════════════════════════════════

\subsection{Class \texttt{FeatureExtractor}}

Tải tệp \texttt{visual\_features\_sample.npz} chứa vector đặc trưng hình ảnh
2048 chiều (trích xuất từ pre-trained CNN, tiêu biểu là ResNet-50).

\begin{itemize}[leftmargin=2em]
    \item Ánh xạ \texttt{article\_id} $\to$ vector $\mathbf{v}_i \in \mathbb{R}^{2048}$.
    \item Xử lý key lookup linh hoạt: hỗ trợ cả dạng chuỗi gốc và zero-padded.
    \item Khi không có feature, trả về vector $\mathbf{0}$.
\end{itemize}

\subsection{Class \texttt{NCF} — Neural Collaborative Filtering}

Kiến trúc kết hợp hai nhánh song song:

\begin{enumerate}[leftmargin=2em]
    \item \textbf{GMF (Generalized Matrix Factorization)}:
    \begin{equation}
        \mathbf{h}_{GMF} = \mathbf{p}_u^{MF} \odot \mathbf{q}_i^{MF}
    \end{equation}
    
    \item \textbf{MLP (Multi-Layer Perceptron)}:
    \begin{equation}
        \mathbf{h}_{MLP} = \sigma\!\left(
            \mathbf{W}_L \cdots \sigma\!\left(
                \mathbf{W}_1 \begin{bmatrix} \mathbf{p}_u^{MLP} \\ \mathbf{q}_i^{MLP} \end{bmatrix}
                + \mathbf{b}_1
            \right) \cdots + \mathbf{b}_L
        \right)
    \end{equation}
    
    \item \textbf{Kết hợp}:
    \begin{equation}
        \hat{y}_{ui} = \sigma\!\left(
            \mathbf{w}^{\top} \begin{bmatrix} \mathbf{h}_{GMF} \\ \mathbf{h}_{MLP} \end{bmatrix}
        \right)
    \end{equation}
\end{enumerate}

Cấu hình mặc định: \texttt{mf\_dim=8}, MLP layers: $[64, 32, 16, 8]$,
Dropout = 0.1 sau mỗi lớp ReLU.

\subsection{Class \texttt{HybridModel} — Mô hình Hybrid}

Ghép (Concatenate) vector latent từ NCF với vector visual features,
rồi đưa qua một chuỗi Dense layers:

\begin{equation}
    \hat{y}_{ui}^{hybrid} = \sigma\!\left(
        \text{Dense}\!\left(
            \begin{bmatrix}
                \mathbf{h}_{NCF} \\ \mathbf{v}_i
            \end{bmatrix}
        \right)
    \right)
\end{equation}

Kiến trúc Dense layers:
\begin{center}
\texttt{(ncf\_dim + 2048)} $\to$ \texttt{256} $\to$ \texttt{BN + ReLU + Dropout(0.2)}
$\to$ \texttt{64} $\to$ \texttt{BN + ReLU + Dropout(0.1)}
$\to$ \texttt{1} $\to$ \texttt{Sigmoid}
\end{center}

% ═════════════════════════════════════════════════════════════════════════════
\section{Module \texttt{metrics.py} — Đánh giá hiệu năng}
% ═════════════════════════════════════════════════════════════════════════════

Module tập trung các hàm tính toán chỉ số đánh giá:

\subsection{MAP@12 — Mean Average Precision at 12}

Chỉ số chính theo cuộc thi Kaggle H\&M:

\begin{equation}
    \text{MAP@12} = \frac{1}{|U|} \sum_{u \in U}
    \frac{1}{\min(m_u, 12)}
    \sum_{k=1}^{12} P(k) \cdot \text{rel}(k)
\end{equation}

trong đó:
\begin{itemize}[leftmargin=2em]
    \item $|U|$ là tổng số người dùng trong tập test.
    \item $m_u$ là số sản phẩm thực tế mà người dùng $u$ đã mua.
    \item $P(k)$ là precision tại vị trí $k$.
    \item $\text{rel}(k) = 1$ nếu sản phẩm tại vị trí $k$ nằm trong tập thực tế.
\end{itemize}

\subsection{Các chỉ số bổ sung}

\begin{itemize}[leftmargin=2em]
    \item \textbf{Hit Rate@K}: Tỉ lệ người dùng có ít nhất 1 gợi ý đúng trong top-K.
    \item \textbf{NDCG@K}: Normalized Discounted Cumulative Gain —
          cân nhắc vị trí xếp hạng của các gợi ý đúng.
\end{itemize}

% ═════════════════════════════════════════════════════════════════════════════
\section{Script \texttt{train\_hybrid.py} — Huấn luyện mô hình}
% ═════════════════════════════════════════════════════════════════════════════

\subsection{Negative Sampling}

Với mỗi cặp tương tác thực $(u, i^+)$, sinh ra $N$ mẫu âm $(u, i^-)$
bằng cách chọn ngẫu nhiên các sản phẩm mà người dùng \textit{chưa} mua:

\begin{equation}
    \mathcal{D} = \{(u, i^+, y=1)\} \cup \{(u, i^-, y=0) \mid i^- \notin \mathcal{I}_u\}
\end{equation}

Negative samples được \textbf{tái sinh mỗi epoch} để tăng tính đa dạng.

\subsection{Mixed Precision Training}

Sử dụng \texttt{torch.autocast} và \texttt{GradScaler} để huấn luyện
ở chế độ FP16 trên GPU:

\begin{itemize}[leftmargin=2em]
    \item Giảm $\sim$50\% bộ nhớ GPU.
    \item Tăng tốc $\sim$2x trên GPU hỗ trợ Tensor Cores (V100, A100, RTX 30xx+).
    \item Gradient clipping (\texttt{max\_norm=1.0}) để ngăn gradient explosion.
\end{itemize}

\subsection{Early Stopping}

\begin{itemize}[leftmargin=2em]
    \item Theo dõi \texttt{val\_loss} sau mỗi epoch.
    \item Nếu không cải thiện quá \texttt{min\_delta} trong \texttt{patience} epochs
          liên tiếp $\to$ dừng huấn luyện.
    \item Chỉ lưu checkpoint khi đạt \texttt{val\_loss} tốt nhất mới.
\end{itemize}

\subsection{Checkpoint Strategy}

\begin{table}[H]
\centering
\begin{tabular}{ll}
\toprule
\textbf{Tệp} & \textbf{Nội dung} \\
\midrule
\texttt{checkpoints/hybrid\_best.pt} & Weights tốt nhất (val\_loss thấp nhất) \\
\texttt{checkpoints/hybrid\_last.pt} & Full state: model + optimizer + scheduler + scaler \\
\bottomrule
\end{tabular}
\end{table}

% ═════════════════════════════════════════════════════════════════════════════
\section{Module \texttt{app.py} — Giao diện Streamlit}
% ═════════════════════════════════════════════════════════════════════════════

Ứng dụng web tương tác cho phép:

\begin{enumerate}[leftmargin=2em]
    \item \textbf{Chọn khách hàng}: Chọn từ danh sách hoặc nhập trực tiếp \texttt{user\_id}
          (dữ liệu từ \texttt{customers\_fixed.parquet}).
    \item \textbf{Chọn phương pháp gợi ý}: Hybrid Model hoặc Popularity Baseline.
    \item \textbf{Hiển thị hồ sơ}: Thông tin khách hàng (tuổi, trạng thái thành viên,
          tần suất nhận tin thời trang).
    \item \textbf{Hiển thị Top-12 gợi ý}: Mỗi sản phẩm hiển thị dưới dạng thẻ
          (card) gồm tên sản phẩm, loại, màu sắc, bộ phận, nhóm sản phẩm và
          điểm dự đoán (prediction score).
    \item \textbf{Fallback thông minh}: Tự động chuyển sang Popularity Baseline
          nếu không tìm thấy tệp checkpoint.
\end{enumerate}

\subsection{Luồng hoạt động}

\begin{enumerate}[leftmargin=2em]
    \item Người dùng chọn \texttt{customer\_id} từ sidebar.
    \item Nhấn nút ``Get Recommendations''.
    \item Hệ thống tải model từ \texttt{checkpoints/hybrid\_best.pt}.
    \item Score tất cả sản phẩm ứng viên cho người dùng đã chọn.
    \item Sắp xếp theo score giảm dần, lấy Top-12.
    \item Join metadata từ \texttt{articles\_cleaned.parquet} và hiển thị.
\end{enumerate}

% ═════════════════════════════════════════════════════════════════════════════
\section{Sơ đồ phụ thuộc giữa các module}
% ═════════════════════════════════════════════════════════════════════════════

\begin{figure}[H]
\centering
\begin{verbatim}
    ┌──────────────┐     ┌──────────────┐
    │ data_loader  │     │   models     │
    │   .py        │     │   .py        │
    │              │     │              │
    │ DataLoader   │     │ NCF          │
    │ Polars       │     │ FeatureExtr. │
    │ Popularity   │     │ HybridModel  │
    │ Baseline     │     │              │
    │ MF           │     │              │
    │ calculate_   │     │              │
    │ map12        │     │              │
    └──────┬───────┘     └──────┬───────┘
           │                    │
           ▼                    ▼
    ┌──────────────────────────────────┐
    │       train_hybrid.py            │
    │                                  │
    │  NegativeSamplingDataset         │
    │  EarlyStopping                   │
    │  Training Loop + Mixed Precision │
    └──────────────┬───────────────────┘
                   │
                   ▼
          checkpoints/hybrid_best.pt
                   │
                   ▼
    ┌──────────────────────────────────┐
    │          app.py                  │
    │                                  │
    │  Streamlit UI                    │
    │  Load Model → Score → Display   │
    └──────────────────────────────────┘

    ┌──────────────┐
    │ metrics.py   │  ← standalone evaluation utilities
    │              │
    │ MAP@12       │
    │ Hit Rate     │
    │ NDCG@K       │
    └──────────────┘
\end{verbatim}
\caption{Sơ đồ phụ thuộc giữa các module trong hệ thống}
\end{figure}

% ═════════════════════════════════════════════════════════════════════════════
\section{Hướng dẫn chạy}
% ═════════════════════════════════════════════════════════════════════════════

\subsection{Cài đặt thư viện}
\begin{lstlisting}[language=bash]
pip install polars torch numpy streamlit
\end{lstlisting}

\subsection{Huấn luyện mô hình}
\begin{lstlisting}[language=bash]
python train_hybrid.py --data-dir Dataset --num-epochs 30 --batch-size 4096
\end{lstlisting}

\subsection{Chạy giao diện demo}
\begin{lstlisting}[language=bash]
streamlit run app.py
\end{lstlisting}

% ═════════════════════════════════════════════════════════════════════════════
\section{Kết luận}
% ═════════════════════════════════════════════════════════════════════════════

Hệ thống đã được xây dựng theo kiến trúc module hoá rõ ràng, tách biệt
giữa phần xử lý dữ liệu (\texttt{data\_loader.py}), định nghĩa mô hình
(\texttt{models.py}), đánh giá (\texttt{metrics.py}), huấn luyện
(\texttt{train\_hybrid.py}), và giao diện người dùng (\texttt{app.py}).
Cách tiếp cận Hybrid cho phép tận dụng đồng thời thông tin hành vi mua sắm
(collaborative signals) và đặc trưng thị giác (visual features) của sản phẩm,
hứa hẹn cải thiện đáng kể chất lượng gợi ý so với baseline đơn thuần.

\end{document}
